\section{\texorpdfstring{$\R$}{R}的连通集}

\begin{theorem}{$\R$中的区间=$\R$中的连通集}{}
	设$A\subseteq\R$,则$A$是$\R$的连通集$\iff$ $A$是$\R$的区间.
\end{theorem}

\begin{remark}{}{}
	由此可见,如果$\R$中每个基数$\geq 1$的区间都不被$A$包含,那么$A$在$\R$中完全不连通.例如,$\Q$在$\R$中完全不连通.
\end{remark}

\begin{corollary}{$\R$的连通性和局部连通性}{}
	实直线既连通又局部连通.
\end{corollary}

\begin{corollary}{$\R$的连通紧集都是区间}{}
	$\R$的连通紧集有且仅有$\R$中的有界闭区间.
\end{corollary}

\begin{proposition}{$\R$中的开集分解定理}{}
	对$\R$的每个非空开集$U$,存在唯一的至多可数集$I$和$\R$中一族两两不交的开区间$\left(U_{i}\right)_{i\in I}$使得$U=\bigcup\limits_{i\in I}U_{i}$.
\end{proposition}

\begin{corollary}{$\R$中的闭集分解定理}{}
	\begin{enumerate}
		\item 对$\R$的每非空闭集$F$,存在唯一的至多可数集$I$和$\R$中一族两两不交的开区间$\left(U_{i}\right)_{i\in I}$使得$F=\R\setminus\bigcup\limits_{i\in I}U_{i}$.
		\item 对$\R$中每一族两两不交的开区间$\left(U_{i}\right)_{i\in I}$,集合$\R\setminus\bigcup\limits_{i\in I}U_{i}$是$\R$的闭集.
	\end{enumerate}
\end{corollary}

\begin{definition}{相邻区间}{}
	设$F$是$\R$的闭集,$\left(U_{i}\right)_{i\in I}$是$\R$中两两不交的开区间族.若$U=\R\setminus\bigcup\limits_{i\in I}U_{i}$,则称$\left(U_{i}\right)_{i\in I}$的各项为$F$的构成区间.
\end{definition}

\begin{example}{Cantor三分集的构造}{}
	设$A=\left[0,1\right]\subseteq\R$.对每个$n\geq 0$和$1\leq p\leq 2^{n}$,令序列$\left(b_{i}\right)_{i=1}^{n}$是$p-1$的$2$进制展开式,
	\begin{align*}
		I_{n,p}=\left(\frac{1}{3^{n+1}},\frac{2}{3^{n+1}}\right)+\sum\limits_{i=1}^{n}\frac{2b_{i}}{3^{i}},K_{n,p}=\left[0,\frac{1}{3^{n+1}}\right]+\sum\limits_{i=1}^{n+1}\frac{2b_{i}}{3^{i}}.
	\end{align*}
	易见$\left[0,1\right]\setminus\bigcup\limits_{n=0}^{\infty}\bigcup\limits_{p=1}^{2^{n}}I_{n,p}=\bigcap\limits_{n=0}^{\infty}\bigcup\limits_{p=1}^{2^{n+1}}K_{n,p}$.我们称$\left[0,1\right]\setminus\bigcup\limits_{n=0}^{\infty}\bigcup\limits_{p=1}^{2^{n}}I_{n,p}$是Cantor三分集,它是$\R$中完全不连通的紧集.
\end{example}